% =====================================================================
\section{Evaluation}\label{sec:evaluation}

\subsection{Research Questions}
The guarantees of Section~\ref{sec:guarantees} are proved on paper. The evaluation asks whether they hold in executable instantiations, whether each mechanism is necessary, what exactness buys, and what happens with real language models. RQ1 and RQ3 are \emph{verification} activities (mutation-based guard-necessity checks and conformance tests), not empirical claims about multi-agent practice; RQ5 and RQ6 provide the latter.
\begin{itemize}
\item[\RQ{1}] \emph{Guard necessity.} Does the controller preserve its safety properties under adversarial executions, and does removing any single mechanism break the property it is meant to protect?
\item[\RQ{2}] \emph{Locality, cost and scheduling.} How much generation does exact dependency tracking avoid, on independent and on coupled workspaces, and how does the upstream gate affect repair cost under asynchronous scheduling?
\item[\RQ{3}] \emph{Conformance.} Does the reference implementation handle each class of workflow change exactly as the semantics prescribes?
\item[\RQ{4}] \emph{Assurance.} What do the mechanized proofs and bounded model checks establish, and what remains open?
\item[\RQ{5}] \emph{Real LLM proposers.} With real models generating designs, code and tests, how often are proposals rejected, repaired, globally inconsistent, or stale, how often does early cutoff apply, and what does exactness save?
\item[\RQ{6}] \emph{Validator strength on standard benchmarks.} On HumanEval+ and MBPP+, how often does a value accepted by the declared validator fail stronger tests, and how often are regenerated values equal?
\end{itemize}

\subsection{Experimental Setup}
\emph{Subjects.} RQ1--RQ4 use the scripted case study (Section~\ref{sec:case-study}), the generic kernel and the TLA$^+$ model. RQ5 and RQ6 use the LLM workflows of Section~\ref{sec:llm-workflows} with the generic kernel.

\emph{Scripted and real proposers.} For RQ1--RQ3 the proposal source is a script, for two reasons. The safety guarantees quantify over \emph{all} proposals, so an adversarial script that can return anything is the stronger test. And deterministic proposals make every run reproducible from its seed. The adversarial provider returns valid proposals, locally valid but structurally or jointly inconsistent proposals, malformed values, and plausible answers to a \emph{different} policy question, and it delivers replies late, out of order, duplicated or never. Scripts cannot answer how often real models produce such proposals, whether early cutoff ever applies to stochastic outputs, or what the savings are worth in tokens and time. RQ5 and RQ6 therefore use real models. They run locally through Ollama~\cite{ollama} with fixed model digests, per-request seeds and logged responses: \texttt{qwen2.5-coder:7b}~\cite{hui2024qwen25coder}, \texttt{llama3.1:8b}~\cite{grattafiori2024llama3}, \texttt{gemma2:9b}~\cite{gemma2024gemma2}, \texttt{phi4:14b}~\cite{abdin2024phi4} and \texttt{mistral-small:24b} on NVIDIA H100 and A100 GPUs of the Digital Research Alliance of Canada, and additionally \texttt{qwen3-coder:30b}~\cite{yang2025qwen3} and \texttt{gpt-oss:20b}~\cite{openai2025gptoss} on a local NVIDIA GB10 for RQ5.

\emph{RQ1 ablation harness.} We drive the reference controller and ten variants that each remove exactly one mechanism (two remove both obsolete-reply defenses, as in the formal mutants). An eleventh variant restores the original semantics with structural checks for enabled consumers only (Section~\ref{sec:deferred}). Every variant reuses the reference code and overrides only the affected check. Each run starts with one to three documents, budget $k=2$, and executes \AblSteps{} random actions: dispatch; delivery in random order, possibly duplicated; timeout; source edits that flip a policy, add or delete a document; context updates, including equal-content rewrites, unselected findings and availability toggles; contract changes; owner revocation and regrant; adding or removing the reviewer; cooperative repair bursts; and publication attempts. All variants use the same \AblRuns{} seeds. An \emph{independent oracle} judges the state after every transaction. It recomputes each job's observation from the authoritative inputs, remembers the observation each request was dispatched against, tracks episodes itself, and never reads the controller's keys. It checks seven safety properties and one liveness property:
P1 no value is usable whose evidence was produced for a different or unavailable observation (and no request is sent on one);
P2 no usable value is invalid for its current observation;
P3 no value is usable while a prerequisite is unusable;
P4 no value from a revoked owner is usable;
P5 no ill-typed structure is installed for an active consumer;
P6 no episode receives more than $k$ requests;
P7 every published snapshot is fully usable and satisfies $G$;
L1 no authorized edit or context change is refused because of a committed generated value.
The violation rates are a property of the action generator's sampling distribution, not of any deployment. What matters is whether a property is violated at all.

\emph{RQ2 harnesses.} Five strategies share validators, scheduler and guards and differ \emph{only} in what a job's key observes: (i) \emph{Exact} (this paper); (ii) \emph{No early cutoff}, where a job also depends on the identity of each upstream acceptance, as in timestamp-based build tools; (iii) \emph{Version-keyed context}, where context readers depend on the global context version; (iv) \emph{Whole workspace}, where every job sees all inputs and the whole context, as in shared-message-pool frameworks; and (v) \emph{Full regeneration}. The \emph{independent} workspace replicates the four-view chain for $N\in\{1,2,4,\ldots,\CostMaxN\}$ mutually independent documents, with \CostSeeds{} seeded streams of \CostChanges{} changes per size. The \emph{coupled} workspace has \CoupledModules{} modules whose design jobs read $f\in\{0,1,2,4,8\}$ of \CoupledShared{} shared data-model components. Each component is a two-job chain, a specification with fields and a description followed by a derived schema of the fields. Its change stream adds semantic edits (a new field) and cosmetic edits (description only) of shared components. The \emph{scheduling} harness uses layered lattices of depth $d\le\SchedMaxDepth$ and width $w\le\SchedMaxWidth$ in which every node reads all nodes of the previous layer. After one root edit it repairs under random request latencies with unlimited parallelism, with and without the upstream gate (\SchedSeeds{} seeds per cell). Providers are deterministic functions of the observation, so all strategies must accept identical values; we check this after every change.

\emph{RQ3, RQ4.} The \Scenarios{} scenarios and their expectations are fixed in advance in \texttt{campaign.json}. Model checking uses TLC (version \TLCExtVersion). Proofs use TLAPS with its bundled Z3 4.8.9 in strict mode with fingerprints erased. Python \PythonVersion{} was used for the controllers.

\emph{RQ5 protocol.} For each model, we run \emph{completion} trials: a fresh four-service workspace with random policies and findings is repaired from scratch with $k=3$ and validator feedback. If the barrier reports a conflict, one feedback epoch promotes the majority identifier type to a workspace convention. In \emph{stale} trials, an implementation request is dispatched, R03 is flipped half a second later while the model is still generating, and the late reply is delivered. In \emph{cutoff} trials, at temperatures 0 and 0.7, the design contract and then the implementation contract are replaced by a reworded instruction with the same meaning, and we compare the regenerated value with the previous one. In \emph{cost} runs, one fixed stream of eight changes is applied under four strategies. All other generation settings are fixed (temperature 0.7, at most 2,500 output tokens).

\emph{RQ6 protocol.} Each model runs all \BenchTasks{} tasks once at temperature 0.2 with $k=3$. After the initial repair, every code job's contract is replaced by a reworded instruction with the same meaning, and the workflow is repaired again.

% ---------------------------------------------------------------------
\subsection{RQ1: Necessity of Each Guard}\label{sec:rq1}
\emph{Bounded model checking.} The original nine baselines, nine one-job mutants and five nonvacuity witnesses were re-run for this revision with the current TLC release, and every state count is identical to the original submission (supplement, Table~S1); each one-job mutant violated its targeted property. Table~\ref{tab:tlc-ext} adds multi-job instances (RQ4). There, each of the \TLCExtMutants{} mutant configurations, every mutant on chains of two and three jobs and on a four-job diamond, behaves as predicted: \TLCExtMutantsCaught{} violate their targeted property within a few actions, and only the six configurations that remove a single obsolete-reply defense survive, because the other defense still blocks the stale commit.


\emph{Implementation-level ablation.} Fig.~\ref{fig:ablation} asks the same question of the implementation, with several documents, concurrent replies and adversarial proposals. Across \AblRuns{} runs, the reference controller showed \textbf{no} violation of any property (\FullViolRuns{} runs) and refused no authorized change (\FullRefusals{} refusals). It still published in \FullPubRuns{} runs (\FullPublications{} publications), so the absence of violations is not due to inactivity. Every non-redundant ablation was detected on the property it targets:
\begin{itemize}
\item removing both obsolete-reply defenses exposed stale values in \ObsoleteStalePct\% of runs, and in \ObsoleteInvalidPct\% values that were \emph{invalid} for the current input;
\item omitting the contract from the key did so in \ContractStalePct\%, and omitting selected context in \ContextStalePct\%;
\item removing the upstream gate exposed values over stale prerequisites in \UpstreamPct\%;
\item ignoring authority exposed revoked owners' values in \AuthorityPct\%;
\item skipping candidate normalization installed ill-typed structure in \CandidatePct\%;
\item resetting budgets exceeded $k$ requests per episode in \BudgetPct\%;
\item dropping the global predicate published inconsistent snapshots in \GlobalPublishPct\%;
\item checking structural expressions only for enabled consumers refused \EnabledOnlyRefusals{} authorized changes in \EnabledOnlyRuns{} runs (L1), without any safety violation.
\end{itemize}
Removing \emph{one} obsolete-reply defense again produced no violation, which replicates the model-checking result at implementation scale. Several ablations also violate secondary properties; stale contracts, for example, propagate into P3 and P7, because a stale value feeds its descendants and can reach publication.

% Generated by experiments/tse_paper.py; do not edit.
\begin{figure}[!t]
\centering
\begin{tikzpicture}
\node[rotate=60,anchor=south west,font=\scriptsize] at (-0.12,0.02) {P1 Stale};
\node[rotate=60,anchor=south west,font=\scriptsize] at (0.44,0.02) {P2 Invalid};
\node[rotate=60,anchor=south west,font=\scriptsize] at (1.00,0.02) {P3 Upstream};
\node[rotate=60,anchor=south west,font=\scriptsize] at (1.56,0.02) {P4 Authority};
\node[rotate=60,anchor=south west,font=\scriptsize] at (2.12,0.02) {P5 Structure};
\node[rotate=60,anchor=south west,font=\scriptsize] at (2.68,0.02) {P6 Budget};
\node[rotate=60,anchor=south west,font=\scriptsize] at (3.24,0.02) {P7 Publish};
\node[rotate=60,anchor=south west,font=\scriptsize] at (3.80,0.02) {L1 Refusal};
\node[rotate=60,anchor=south west,font=\scriptsize] at (4.36,0.02) {Published};
\node[anchor=east,font=\scriptsize] at (-0.33,-0.20) {None (reference)};
\fill[resultteal!12!white] (-0.26,-0.38) rectangle ++(0.52,0.36);
\node[font=\tiny,text=resultteal!70!black] at (0.00,-0.20) {0};
\fill[resultteal!12!white] (0.30,-0.38) rectangle ++(0.52,0.36);
\node[font=\tiny,text=resultteal!70!black] at (0.56,-0.20) {0};
\fill[resultteal!12!white] (0.86,-0.38) rectangle ++(0.52,0.36);
\node[font=\tiny,text=resultteal!70!black] at (1.12,-0.20) {0};
\fill[resultteal!12!white] (1.42,-0.38) rectangle ++(0.52,0.36);
\node[font=\tiny,text=resultteal!70!black] at (1.68,-0.20) {0};
\fill[resultteal!12!white] (1.98,-0.38) rectangle ++(0.52,0.36);
\node[font=\tiny,text=resultteal!70!black] at (2.24,-0.20) {0};
\fill[resultteal!12!white] (2.54,-0.38) rectangle ++(0.52,0.36);
\node[font=\tiny,text=resultteal!70!black] at (2.80,-0.20) {0};
\fill[resultteal!12!white] (3.10,-0.38) rectangle ++(0.52,0.36);
\node[font=\tiny,text=resultteal!70!black] at (3.36,-0.20) {0};
\fill[resultteal!12!white] (3.66,-0.38) rectangle ++(0.52,0.36);
\node[font=\tiny,text=resultteal!70!black] at (3.92,-0.20) {0};
\fill[gray!32!white] (4.22,-0.38) rectangle ++(0.52,0.36);
\node[font=\tiny,text=black] at (4.48,-0.20) {69};
\draw[inkblue!60,densely dotted] (-3.3,-0.40) -- (4.76,-0.40);
\node[anchor=east,font=\scriptsize] at (-0.33,-0.60) {Cancellation};
\fill[resultteal!12!white] (-0.26,-0.78) rectangle ++(0.52,0.36);
\node[font=\tiny,text=resultteal!70!black] at (0.00,-0.60) {0};
\fill[resultteal!12!white] (0.30,-0.78) rectangle ++(0.52,0.36);
\node[font=\tiny,text=resultteal!70!black] at (0.56,-0.60) {0};
\fill[resultteal!12!white] (0.86,-0.78) rectangle ++(0.52,0.36);
\node[font=\tiny,text=resultteal!70!black] at (1.12,-0.60) {0};
\fill[resultteal!12!white] (1.42,-0.78) rectangle ++(0.52,0.36);
\node[font=\tiny,text=resultteal!70!black] at (1.68,-0.60) {0};
\fill[resultteal!12!white] (1.98,-0.78) rectangle ++(0.52,0.36);
\node[font=\tiny,text=resultteal!70!black] at (2.24,-0.60) {0};
\fill[resultteal!12!white] (2.54,-0.78) rectangle ++(0.52,0.36);
\node[font=\tiny,text=resultteal!70!black] at (2.80,-0.60) {0};
\fill[resultteal!12!white] (3.10,-0.78) rectangle ++(0.52,0.36);
\node[font=\tiny,text=resultteal!70!black] at (3.36,-0.60) {0};
\fill[resultteal!12!white] (3.66,-0.78) rectangle ++(0.52,0.36);
\node[font=\tiny,text=resultteal!70!black] at (3.92,-0.60) {0};
\fill[gray!31!white] (4.22,-0.78) rectangle ++(0.52,0.36);
\node[font=\tiny,text=black] at (4.48,-0.60) {68};
\node[anchor=east,font=\scriptsize] at (-0.33,-1.00) {Commit recheck};
\fill[resultteal!12!white] (-0.26,-1.18) rectangle ++(0.52,0.36);
\node[font=\tiny,text=resultteal!70!black] at (0.00,-1.00) {0};
\fill[resultteal!12!white] (0.30,-1.18) rectangle ++(0.52,0.36);
\node[font=\tiny,text=resultteal!70!black] at (0.56,-1.00) {0};
\fill[resultteal!12!white] (0.86,-1.18) rectangle ++(0.52,0.36);
\node[font=\tiny,text=resultteal!70!black] at (1.12,-1.00) {0};
\fill[resultteal!12!white] (1.42,-1.18) rectangle ++(0.52,0.36);
\node[font=\tiny,text=resultteal!70!black] at (1.68,-1.00) {0};
\fill[resultteal!12!white] (1.98,-1.18) rectangle ++(0.52,0.36);
\node[font=\tiny,text=resultteal!70!black] at (2.24,-1.00) {0};
\fill[resultteal!12!white] (2.54,-1.18) rectangle ++(0.52,0.36);
\node[font=\tiny,text=resultteal!70!black] at (2.80,-1.00) {0};
\fill[resultteal!12!white] (3.10,-1.18) rectangle ++(0.52,0.36);
\node[font=\tiny,text=resultteal!70!black] at (3.36,-1.00) {0};
\fill[resultteal!12!white] (3.66,-1.18) rectangle ++(0.52,0.36);
\node[font=\tiny,text=resultteal!70!black] at (3.92,-1.00) {0};
\fill[gray!32!white] (4.22,-1.18) rectangle ++(0.52,0.36);
\node[font=\tiny,text=black] at (4.48,-1.00) {69};
\node[anchor=east,font=\scriptsize] at (-0.33,-1.40) {Cancellation + recheck};
\fill[resultamber!51!white] (-0.26,-1.58) rectangle ++(0.52,0.36);
\node[font=\tiny,text=black] at (0.00,-1.40) {35};
\fill[resultamber!43!white] (0.30,-1.58) rectangle ++(0.52,0.36);
\node[font=\tiny,text=black] at (0.56,-1.40) {25};
\fill[resultamber!44!white] (0.86,-1.58) rectangle ++(0.52,0.36);
\node[font=\tiny,text=black] at (1.12,-1.40) {26};
\fill[resultteal!12!white] (1.42,-1.58) rectangle ++(0.52,0.36);
\node[font=\tiny,text=resultteal!70!black] at (1.68,-1.40) {0};
\fill[resultteal!12!white] (1.98,-1.58) rectangle ++(0.52,0.36);
\node[font=\tiny,text=resultteal!70!black] at (2.24,-1.40) {0};
\fill[resultteal!12!white] (2.54,-1.58) rectangle ++(0.52,0.36);
\node[font=\tiny,text=resultteal!70!black] at (2.80,-1.40) {0};
\fill[resultamber!34!white] (3.10,-1.58) rectangle ++(0.52,0.36);
\node[font=\tiny,text=black] at (3.36,-1.40) {12};
\fill[resultteal!12!white] (3.66,-1.58) rectangle ++(0.52,0.36);
\node[font=\tiny,text=resultteal!70!black] at (3.92,-1.40) {0};
\fill[gray!32!white] (4.22,-1.58) rectangle ++(0.52,0.36);
\node[font=\tiny,text=black] at (4.48,-1.40) {69};
\node[anchor=east,font=\scriptsize] at (-0.33,-1.80) {Contract in signature};
\fill[resultamber!85!white] (-0.26,-1.98) rectangle ++(0.52,0.36);
\node[font=\tiny,text=white] at (0.00,-1.80) {81};
\fill[resultteal!12!white] (0.30,-1.98) rectangle ++(0.52,0.36);
\node[font=\tiny,text=resultteal!70!black] at (0.56,-1.80) {0};
\fill[resultamber!77!white] (0.86,-1.98) rectangle ++(0.52,0.36);
\node[font=\tiny,text=white] at (1.12,-1.80) {70};
\fill[resultteal!12!white] (1.42,-1.98) rectangle ++(0.52,0.36);
\node[font=\tiny,text=resultteal!70!black] at (1.68,-1.80) {0};
\fill[resultteal!12!white] (1.98,-1.98) rectangle ++(0.52,0.36);
\node[font=\tiny,text=resultteal!70!black] at (2.24,-1.80) {0};
\fill[resultteal!12!white] (2.54,-1.98) rectangle ++(0.52,0.36);
\node[font=\tiny,text=resultteal!70!black] at (2.80,-1.80) {0};
\fill[resultamber!47!white] (3.10,-1.98) rectangle ++(0.52,0.36);
\node[font=\tiny,text=black] at (3.36,-1.80) {30};
\fill[resultteal!12!white] (3.66,-1.98) rectangle ++(0.52,0.36);
\node[font=\tiny,text=resultteal!70!black] at (3.92,-1.80) {0};
\fill[gray!31!white] (4.22,-1.98) rectangle ++(0.52,0.36);
\node[font=\tiny,text=black] at (4.48,-1.80) {66};
\node[anchor=east,font=\scriptsize] at (-0.33,-2.20) {Context in footprint};
\fill[resultamber!57!white] (-0.26,-2.38) rectangle ++(0.52,0.36);
\node[font=\tiny,text=black] at (0.00,-2.20) {44};
\fill[resultteal!12!white] (0.30,-2.38) rectangle ++(0.52,0.36);
\node[font=\tiny,text=resultteal!70!black] at (0.56,-2.20) {0};
\fill[resultamber!53!white] (0.86,-2.38) rectangle ++(0.52,0.36);
\node[font=\tiny,text=black] at (1.12,-2.20) {38};
\fill[resultteal!12!white] (1.42,-2.38) rectangle ++(0.52,0.36);
\node[font=\tiny,text=resultteal!70!black] at (1.68,-2.20) {0};
\fill[resultteal!12!white] (1.98,-2.38) rectangle ++(0.52,0.36);
\node[font=\tiny,text=resultteal!70!black] at (2.24,-2.20) {0};
\fill[resultteal!12!white] (2.54,-2.38) rectangle ++(0.52,0.36);
\node[font=\tiny,text=resultteal!70!black] at (2.80,-2.20) {0};
\fill[resultamber!38!white] (3.10,-2.38) rectangle ++(0.52,0.36);
\node[font=\tiny,text=black] at (3.36,-2.20) {18};
\fill[resultteal!12!white] (3.66,-2.38) rectangle ++(0.52,0.36);
\node[font=\tiny,text=resultteal!70!black] at (3.92,-2.20) {0};
\fill[gray!31!white] (4.22,-2.38) rectangle ++(0.52,0.36);
\node[font=\tiny,text=black] at (4.48,-2.20) {68};
\node[anchor=east,font=\scriptsize] at (-0.33,-2.60) {Upstream gate};
\fill[resultamber!92!white] (-0.26,-2.78) rectangle ++(0.52,0.36);
\node[font=\tiny,text=white] at (0.00,-2.60) {89};
\fill[resultteal!12!white] (0.30,-2.78) rectangle ++(0.52,0.36);
\node[font=\tiny,text=resultteal!70!black] at (0.56,-2.60) {0};
\fill[resultamber!99!white] (0.86,-2.78) rectangle ++(0.52,0.36);
\node[font=\tiny,text=white] at (1.12,-2.60) {99};
\fill[resultteal!12!white] (1.42,-2.78) rectangle ++(0.52,0.36);
\node[font=\tiny,text=resultteal!70!black] at (1.68,-2.60) {0};
\fill[resultteal!12!white] (1.98,-2.78) rectangle ++(0.52,0.36);
\node[font=\tiny,text=resultteal!70!black] at (2.24,-2.60) {0};
\fill[resultteal!12!white] (2.54,-2.78) rectangle ++(0.52,0.36);
\node[font=\tiny,text=resultteal!70!black] at (2.80,-2.60) {0};
\fill[resultamber!70!white] (3.10,-2.78) rectangle ++(0.52,0.36);
\node[font=\tiny,text=white] at (3.36,-2.60) {60};
\fill[resultteal!12!white] (3.66,-2.78) rectangle ++(0.52,0.36);
\node[font=\tiny,text=resultteal!70!black] at (3.92,-2.60) {0};
\fill[gray!35!white] (4.22,-2.78) rectangle ++(0.52,0.36);
\node[font=\tiny,text=black] at (4.48,-2.60) {79};
\node[anchor=east,font=\scriptsize] at (-0.33,-3.00) {Authority check};
\fill[resultamber!88!white] (-0.26,-3.18) rectangle ++(0.52,0.36);
\node[font=\tiny,text=white] at (0.00,-3.00) {84};
\fill[resultteal!12!white] (0.30,-3.18) rectangle ++(0.52,0.36);
\node[font=\tiny,text=resultteal!70!black] at (0.56,-3.00) {0};
\fill[resultamber!88!white] (0.86,-3.18) rectangle ++(0.52,0.36);
\node[font=\tiny,text=white] at (1.12,-3.00) {85};
\fill[resultamber!94!white] (1.42,-3.18) rectangle ++(0.52,0.36);
\node[font=\tiny,text=white] at (1.68,-3.00) {93};
\fill[resultteal!12!white] (1.98,-3.18) rectangle ++(0.52,0.36);
\node[font=\tiny,text=resultteal!70!black] at (2.24,-3.00) {0};
\fill[resultteal!12!white] (2.54,-3.18) rectangle ++(0.52,0.36);
\node[font=\tiny,text=resultteal!70!black] at (2.80,-3.00) {0};
\fill[resultamber!65!white] (3.10,-3.18) rectangle ++(0.52,0.36);
\node[font=\tiny,text=white] at (3.36,-3.00) {54};
\fill[resultteal!12!white] (3.66,-3.18) rectangle ++(0.52,0.36);
\node[font=\tiny,text=resultteal!70!black] at (3.92,-3.00) {0};
\fill[gray!35!white] (4.22,-3.18) rectangle ++(0.52,0.36);
\node[font=\tiny,text=black] at (4.48,-3.00) {79};
\node[anchor=east,font=\scriptsize] at (-0.33,-3.40) {Episode budget};
\fill[resultteal!12!white] (-0.26,-3.58) rectangle ++(0.52,0.36);
\node[font=\tiny,text=resultteal!70!black] at (0.00,-3.40) {0};
\fill[resultteal!12!white] (0.30,-3.58) rectangle ++(0.52,0.36);
\node[font=\tiny,text=resultteal!70!black] at (0.56,-3.40) {0};
\fill[resultteal!12!white] (0.86,-3.58) rectangle ++(0.52,0.36);
\node[font=\tiny,text=resultteal!70!black] at (1.12,-3.40) {0};
\fill[resultteal!12!white] (1.42,-3.58) rectangle ++(0.52,0.36);
\node[font=\tiny,text=resultteal!70!black] at (1.68,-3.40) {0};
\fill[resultteal!12!white] (1.98,-3.58) rectangle ++(0.52,0.36);
\node[font=\tiny,text=resultteal!70!black] at (2.24,-3.40) {0};
\fill[resultamber!46!white] (2.54,-3.58) rectangle ++(0.52,0.36);
\node[font=\tiny,text=black] at (2.80,-3.40) {28};
\fill[resultteal!12!white] (3.10,-3.58) rectangle ++(0.52,0.36);
\node[font=\tiny,text=resultteal!70!black] at (3.36,-3.40) {0};
\fill[resultteal!12!white] (3.66,-3.58) rectangle ++(0.52,0.36);
\node[font=\tiny,text=resultteal!70!black] at (3.92,-3.40) {0};
\fill[gray!33!white] (4.22,-3.58) rectangle ++(0.52,0.36);
\node[font=\tiny,text=black] at (4.48,-3.40) {72};
\node[anchor=east,font=\scriptsize] at (-0.33,-3.80) {Candidate normalization};
\fill[resultteal!12!white] (-0.26,-3.98) rectangle ++(0.52,0.36);
\node[font=\tiny,text=resultteal!70!black] at (0.00,-3.80) {0};
\fill[resultteal!12!white] (0.30,-3.98) rectangle ++(0.52,0.36);
\node[font=\tiny,text=resultteal!70!black] at (0.56,-3.80) {0};
\fill[resultteal!12!white] (0.86,-3.98) rectangle ++(0.52,0.36);
\node[font=\tiny,text=resultteal!70!black] at (1.12,-3.80) {0};
\fill[resultteal!12!white] (1.42,-3.98) rectangle ++(0.52,0.36);
\node[font=\tiny,text=resultteal!70!black] at (1.68,-3.80) {0};
\fill[resultamber!58!white] (1.98,-3.98) rectangle ++(0.52,0.36);
\node[font=\tiny,text=black] at (2.24,-3.80) {44};
\fill[resultteal!12!white] (2.54,-3.98) rectangle ++(0.52,0.36);
\node[font=\tiny,text=resultteal!70!black] at (2.80,-3.80) {0};
\fill[resultamber!28!white] (3.10,-3.98) rectangle ++(0.52,0.36);
\node[font=\tiny,text=black] at (3.36,-3.80) {5};
\fill[resultteal!12!white] (3.66,-3.98) rectangle ++(0.52,0.36);
\node[font=\tiny,text=resultteal!70!black] at (3.92,-3.80) {0};
\fill[gray!30!white] (4.22,-3.98) rectangle ++(0.52,0.36);
\node[font=\tiny,text=black] at (4.48,-3.80) {63};
\node[anchor=east,font=\scriptsize] at (-0.33,-4.20) {Global predicate};
\fill[resultteal!12!white] (-0.26,-4.38) rectangle ++(0.52,0.36);
\node[font=\tiny,text=resultteal!70!black] at (0.00,-4.20) {0};
\fill[resultteal!12!white] (0.30,-4.38) rectangle ++(0.52,0.36);
\node[font=\tiny,text=resultteal!70!black] at (0.56,-4.20) {0};
\fill[resultteal!12!white] (0.86,-4.38) rectangle ++(0.52,0.36);
\node[font=\tiny,text=resultteal!70!black] at (1.12,-4.20) {0};
\fill[resultteal!12!white] (1.42,-4.38) rectangle ++(0.52,0.36);
\node[font=\tiny,text=resultteal!70!black] at (1.68,-4.20) {0};
\fill[resultteal!12!white] (1.98,-4.38) rectangle ++(0.52,0.36);
\node[font=\tiny,text=resultteal!70!black] at (2.24,-4.20) {0};
\fill[resultteal!12!white] (2.54,-4.38) rectangle ++(0.52,0.36);
\node[font=\tiny,text=resultteal!70!black] at (2.80,-4.20) {0};
\fill[resultamber!51!white] (3.10,-4.38) rectangle ++(0.52,0.36);
\node[font=\tiny,text=black] at (3.36,-4.20) {35};
\fill[resultteal!12!white] (3.66,-4.38) rectangle ++(0.52,0.36);
\node[font=\tiny,text=resultteal!70!black] at (3.92,-4.20) {0};
\fill[gray!35!white] (4.22,-4.38) rectangle ++(0.52,0.36);
\node[font=\tiny,text=black] at (4.48,-4.20) {79};
\node[anchor=east,font=\scriptsize] at (-0.33,-4.60) {Configured-consumer checks};
\fill[resultteal!12!white] (-0.26,-4.78) rectangle ++(0.52,0.36);
\node[font=\tiny,text=resultteal!70!black] at (0.00,-4.60) {0};
\fill[resultteal!12!white] (0.30,-4.78) rectangle ++(0.52,0.36);
\node[font=\tiny,text=resultteal!70!black] at (0.56,-4.60) {0};
\fill[resultteal!12!white] (0.86,-4.78) rectangle ++(0.52,0.36);
\node[font=\tiny,text=resultteal!70!black] at (1.12,-4.60) {0};
\fill[resultteal!12!white] (1.42,-4.78) rectangle ++(0.52,0.36);
\node[font=\tiny,text=resultteal!70!black] at (1.68,-4.60) {0};
\fill[resultteal!12!white] (1.98,-4.78) rectangle ++(0.52,0.36);
\node[font=\tiny,text=resultteal!70!black] at (2.24,-4.60) {0};
\fill[resultteal!12!white] (2.54,-4.78) rectangle ++(0.52,0.36);
\node[font=\tiny,text=resultteal!70!black] at (2.80,-4.60) {0};
\fill[resultteal!12!white] (3.10,-4.78) rectangle ++(0.52,0.36);
\node[font=\tiny,text=resultteal!70!black] at (3.36,-4.60) {0};
\fill[resultamber!28!white] (3.66,-4.78) rectangle ++(0.52,0.36);
\node[font=\tiny,text=black] at (3.92,-4.60) {5};
\fill[gray!31!white] (4.22,-4.78) rectangle ++(0.52,0.36);
\node[font=\tiny,text=black] at (4.48,-4.60) {68};
\draw[inkblue!40] (3.64,0.02) -- (3.64,-4.80);
\draw[inkblue!40] (4.20,0.02) -- (4.20,-4.80);
\node[anchor=south east,font=\scriptsize\bfseries] at (-0.33,0.02) {Mechanism removed};
\end{tikzpicture}
\caption{RQ1: implementation-level guard ablation under 300 seeded adversarial executions per variant (120 steps each). Cells give the percentage of runs in which the independent oracle observed at least one violation of property P1--P7 (Section~\ref{sec:rq1}); L1 is the liveness hazard of Section~\ref{sec:deferred} (an authorized change refused because of a committed generated value); the last column gives the percentage of runs that published at least once. The reference controller (top row) shows no violation and no refusal. Removing either obsolete-response defense alone is masked by the other. The bottom row is not a removal but the original semantics, which checked structural expressions only for enabled consumers.}\label{fig:ablation}
\end{figure}


\emph{A defect found by the campaign.} In the original submission, the randomized campaign exposed a defect: \texttt{update\_context} let a normalization exception escape when re-enabled context made a previously committed integer architecture inadmissible (\PreFixApiRuns{} of \PreFixRuns{} runs). We repaired the API at the time. The revision traces the defect to its semantic root, the deferred-check hazard of Section~\ref{sec:deferred}, and repairs the semantics. The commit guard now checks configured consumers, which removes the refusals entirely, and Theorem~\ref{thm:termination} has been corrected accordingly. The L1 column of Fig.~\ref{fig:ablation} quantifies the hazard under the original semantics.

\textbf{Answer to RQ1.} No safety violation occurred in any exhaustive baseline configuration or in the \AblRuns{} adversarial executions of the reference controller. Each ablation of a non-redundant mechanism violated its targeted property, both in the model and in the implementation. The cancellation and recheck defenses are individually redundant but jointly necessary. The revised structural check removes a liveness hazard that the original semantics had.

% ---------------------------------------------------------------------
\subsection{RQ2: Locality, Cost and Scheduling}\label{sec:rq2}
\emph{Independent workspaces.} Fig.~\ref{fig:cost} and Table~\ref{tab:cost} report requests and observation bytes per change for $N$ mutually independent documents. This topology isolates the effect of the key: with exact keys, a source edit regenerates that document's four jobs, a selected-context change regenerates the two jobs that read the finding, and unselected changes and equal-content rewrites regenerate nothing. Only contract changes touch every job of the changed stage. Because the documents are independent, the ratios to the coarser strategies grow with $N$ \emph{by construction}. At $N=\CostMaxN$, exact keys needed \ExactReqMaxN{} requests per change against \FullReqMaxN{} for full regeneration (\RatioFull$\times$), and the whole-workspace strategy sent \BytesRatioWorkspace$\times$ the observation bytes (\WorkspaceKBMaxN{} vs.\ \ExactKBMaxN{} kB per change). These numbers show what exactness is worth when work is independent. They do not predict savings on coupled software, which the next experiment addresses.

% Generated by experiments/tse_paper.py; do not edit.
\begin{figure*}[!t]
\centering
\begin{tikzpicture}
\begin{groupplot}[group style={group size=2 by 1,horizontal sep=1.7cm},vgcplot,width=.43\textwidth,height=5.2cm,
  xmode=log,log basis x=2,ymode=log,xtick={1,2,4,8,16,32},xticklabels={1,2,4,8,16,32},log ticks with fixed point,
  ytick={0.1,0.2,0.5,1,2,5,10,20,50,100,200,500,1000,2000},xlabel={Documents in the workspace ($N$)},
  legend style={at={(1.08,1.04)},anchor=south,legend columns=5,draw=none,/tikz/every even column/.append style={column sep=6pt}}]
\nextgroupplot[ylabel={Generation requests per change}]
\addplot[inkblue,mark=*,thick,mark size=1.6pt] coordinates {(1,1.845) (2,2.34) (4,2.81) (8,4.16) (16,5.62) (32,10.77)};
\addplot[resultteal,mark=square*,thick,mark size=1.6pt] coordinates {(1,2.255) (2,3.08) (4,4.63) (8,7.68) (16,10.5) (32,25.01)};
\addplot[resultamber,mark=triangle*,thick,mark size=1.6pt] coordinates {(1,2.305) (2,3.41) (4,5.48) (8,9.46) (16,17.4) (32,32.24)};
\addplot[violet!70!black,mark=diamond*,thick,mark size=1.6pt] coordinates {(1,2.545) (2,5.51) (4,10.54) (8,20.44) (16,45.6) (32,87.68)};
\addplot[gray!70!black,mark=x,dashed,thick,mark size=1.6pt] coordinates {(1,4) (2,8) (4,16) (8,32) (16,64) (32,128)};
\legend{Exact (ours),No early cutoff,Version-keyed context,Whole workspace,Full regeneration}
\nextgroupplot[ylabel={Observation kB sent per change}]
\addplot[inkblue,mark=*,thick,mark size=1.6pt] coordinates {(1,0.6504) (2,0.8106) (4,0.9113) (8,1.382) (16,1.84) (32,3.404)};
\addplot[resultteal,mark=square*,thick,mark size=1.6pt] coordinates {(1,0.5473) (2,0.7381) (4,1.11) (8,1.87) (16,2.494) (32,5.888)};
\addplot[resultamber,mark=triangle*,thick,mark size=1.6pt] coordinates {(1,0.8893) (2,1.34) (4,2.186) (8,3.899) (16,7.299) (32,13.35)};
\addplot[violet!70!black,mark=diamond*,thick,mark size=1.6pt] coordinates {(1,1.235) (2,2.998) (4,6.741) (8,17.93) (16,53.89) (32,177.7)};
\addplot[gray!70!black,mark=x,dashed,thick,mark size=1.6pt] coordinates {(1,0.9932) (2,1.974) (4,3.907) (8,7.842) (16,15.33) (32,30.82)};
\end{groupplot}
\end{tikzpicture}
\caption{RQ2: generation work per input change as the workspace grows (log--log; mean over 5 seeds $\times$ 40 random changes per size). All strategies share validators, scheduler and guards and differ only in what a job's signature observes; they accepted identical values in every run. Left: dispatched requests. Right: total size of the observations sent to workers, a proxy for prompt input tokens.}\label{fig:cost}
\end{figure*}

% Generated by experiments/tse_paper.py; do not edit.
\begin{table}[!t]
\centering
\caption{RQ2: mean generation requests per change by change type for $N=32$ documents. Exact signatures make the cost of source and context changes independent of $N$; contract changes necessarily touch every job of the changed stage. The $N$ documents are mutually independent (no cross-document dependency); Table~\ref{tab:coupled} adds coupling.}\label{tab:cost}
\scriptsize\setlength{\tabcolsep}{2pt}
\begin{tabular*}{\columnwidth}{@{\extracolsep{\fill}}lrrrrr@{}}\toprule
Change type & \rotatebox{70}{Exact (ours)} & \rotatebox{70}{No early cutoff} & \rotatebox{70}{Version-keyed context} & \rotatebox{70}{Whole workspace} & \rotatebox{70}{Full regeneration}\\\midrule
Source clause edit & 4.0 & 4.0 & 4.0 & 128.0 & 128.0\\
Selected finding changed & 2.0 & 2.0 & 64.0 & 128.0 & 128.0\\
Unselected finding changed & 0.0 & 0.0 & 64.0 & 128.0 & 128.0\\
Equal content, new version & 0.0 & 0.0 & 64.0 & 0.0 & 128.0\\
Contract change & 32.0 & 82.9 & 32.0 & 32.0 & 128.0\\
\midrule All changes & 10.8 & 25.0 & 32.2 & 87.7 & 128.0\\
Bytes per request & 316 & 235 & 414 & 2,027 & 241\\
Total kB per change & 3.4 & 5.9 & 13.3 & 177.7 & 30.8\\
\bottomrule\end{tabular*}
\end{table}


\emph{Coupled workspaces.} Table~\ref{tab:coupled} adds coupling through shared data models. Exact keys now pay for exactly the coupled subgraph a change reaches: a semantic edit of a shared model regenerates the model, its schema and, for every dependent module, design, implementation and verification, so its cost grows linearly with the fan-out, from 2 requests at $f=0$ to \CoupledSemanticFmax{} at $f=\CoupledFmax$. A \emph{cosmetic} edit of a shared model, whose schema regenerates to an equal value, costs \CoupledCosmeticExact{} requests with early cutoff and \CoupledCosmeticNoCut{} without. Over the whole stream, exact keys needed \CoupledExactFzero{} requests per change at $f=0$ and \CoupledExactFmax{} at $f=\CoupledFmax$. The coarser strategies needed \CoupledRatioFullFmax$\times$ (full regeneration), \CoupledRatioWorkspaceFmax$\times$ (whole workspace) and \CoupledRatioCutoffFmax$\times$ (no early cutoff) as many requests at the highest coupling. In every run, every request of the exact strategy was caused by a job whose observation had changed. Exact keys do not make coupled changes cheap; they make the cost equal to the size of the affected subgraph.

% Generated by experiments/revision_paper.py; do not edit.
\begin{table}[!t]
\centering
\caption{RQ2 on a coupled workspace: mean generation requests per change for 32 modules sharing 8 data-model components (144 jobs), 5 seeds $\times$ 60 changes per fan-out $f$ (each module's design reads $f$ shared models). Right: an edit of a shared model that adds a field, and one that changes only its description (exact keys / no early cutoff). With exact keys, requests equalled stale episodes in all runs. Per-change-type results are in the supplement (Table~S2).}\label{tab:coupled}
\footnotesize\setlength{\tabcolsep}{3pt}
\begin{tabular*}{\columnwidth}{@{\extracolsep{\fill}}rrrrrrr@{}}\toprule
& \multicolumn{4}{c}{All changes} & \multicolumn{2}{c}{Shared-model edit}\\\cmidrule(lr){2-5}\cmidrule(l){6-7}
$f$ & Exact & No cutoff & Whole ws. & Full & New field & Description\\\midrule
0 & 9.5 & 21.9 & 118.2 & 144.0 & 2.0 & 2.0 / 2.0\\
1 & 10.7 & 21.8 & 120.1 & 144.0 & 14.2 & 2.0 / 12.6\\
2 & 12.5 & 25.8 & 119.4 & 144.0 & 25.6 & 2.0 / 27.6\\
4 & 16.4 & 31.5 & 120.5 & 144.0 & 50.3 & 2.0 / 53.4\\
8 & 22.1 & 42.3 & 121.6 & 144.0 & 98.0 & 2.0 / 98.0\\
\bottomrule\end{tabular*}
\end{table}


\emph{Scheduling.} Table~\ref{tab:sched} tests Theorem~\ref{thm:single} and Proposition~\ref{prop:eager}. With the upstream gate, every affected job started exactly one stale episode and was dispatched once, for every seed and lattice, so requests equal the number of jobs. Without the gate, jobs were dispatched on partially repaired predecessors and re-dispatched as each predecessor committed: up to \SchedEagerReqRatioMax$\times$ the requests and up to \SchedEagerMaxEpisodes{} episodes for a single job, growing with width and depth. The ungated scheduler bought no latency in exchange: makespans differed by at most \SchedMakespanDiffPct\%, because the speculative work was discarded whenever a predecessor changed.

% Generated by experiments/revision_paper.py; do not edit.
\begin{table}[!t]
\centering
\caption{RQ2 scheduling: requests after one root edit on layered lattices (depth $d$, width $w$, every node reads all nodes of the previous layer), mean of 30 seeds with random latencies and unlimited parallelism. Cells: gated requests / ungated requests / worst episodes of one job (ungated). The gated scheduler started exactly one stale episode per affected job in every run (Theorem~\ref{thm:single}); makespans differed by at most 5\%.}\label{tab:sched}
\scriptsize\setlength{\tabcolsep}{3pt}
\begin{tabular*}{\columnwidth}{@{\extracolsep{\fill}}lcccc@{}}\toprule
$d$ & $w=1$ & $w=2$ & $w=3$ & $w=4$\\\midrule
2 & 2 / 2.0 / 1 & 4 / 6.0 / 2 & 6 / 12.0 / 3 & 8 / 20.0 / 4\\
3 & 3 / 3.0 / 1 & 6 / 10.1 / 3 & 9 / 21.3 / 5 & 12 / 36.3 / 5\\
4 & 4 / 4.0 / 1 & 8 / 14.2 / 3 & 12 / 30.8 / 5 & 16 / 52.5 / 6\\
5 & 5 / 5.0 / 1 & 10 / 18.6 / 4 & 15 / 39.9 / 5 & 20 / 69.7 / 7\\
6 & 6 / 6.0 / 1 & 12 / 25.1 / 5 & 18 / 54.3 / 6 & 24 / 85.3 / 7\\
\bottomrule\end{tabular*}
\end{table}


The scenario campaign confirms the mechanism on individual interventions. A repair pass with no relevant change (S02), an unselected context update (S08) and a new context version with equal selected content (S10) caused no request. An architecture contract change repaired to an equal output caused exactly one request, because early cutoff restored the descendants (S16).

\textbf{Answer to RQ2.} Exact keys make the cost of a change equal to the size of the subgraph whose observations it changes. On independent workspaces this is constant in workspace size, and on coupled workspaces it grows with fan-out. Against full regeneration, the saving was \RatioFull$\times$ on \CostMaxN{} independent documents and still \CoupledRatioFullFmax$\times$ at the highest coupling we measured. Early cutoff matters most for changes that do not propagate semantically. The upstream gate makes repair cost independent of scheduling, at no latency cost in our simulations. RQ5 and RQ6 measure what these structural savings mean with stochastic outputs.

% ---------------------------------------------------------------------
\subsection{RQ3: Conformance Under Change}\label{sec:rq3}
Table~\ref{tab:scenarios} groups the \Scenarios{} scenarios by the class of change they exercise. All \ScenarioPasses{} met their predefined expectations. Many expectations are refusals: a discarded late reply (S11--S13), an exhausted episode that earns no new budget (S14), a blocked descendant (S15), a joint conflict (S18), a rolled-back candidate (S19), a reviewer whose addition blocks publication until it has assessed every operation (S20), and a rejected publication based on a stale check (S25). The trace checker admitted all \Scenarios{} recorded traces and rejected all \Corruptions{} deliberately corrupted ones, including acceptance events that hide a source, context or contract change behind a valid value. In S12 the clause changes and then reverts, so the observation returns to its original value. The cancelled request still cannot commit, because the reversion started a new episode: keys alone do not give idempotent delivery, and request identities are needed as well. The expectations were written by the author, so these are conformance tests of the implementation against the semantics, not independent evidence about the semantics itself.

% Generated by experiments/tse_paper.py; do not edit.
\begin{table*}[!t]
\centering
\caption{The 28 scripted change scenarios, grouped by the class of change they exercise. Event counts come from the recorded traces: D dispatched requests, A accepted replies, R current rejections or timeouts, X obsolete or duplicate replies discarded, E committed edits/reconfigurations, B refused edits or publications, P publications. A scenario passes when its independently written expectation holds; for several scenarios the expected outcome is a refusal.}\label{tab:scenarios}
\footnotesize\setlength{\tabcolsep}{4pt}
\begin{tabular*}{\textwidth}{@{\extracolsep{\fill}}llcrrrrrrrc@{}}\toprule
Change class & Scenarios (S\#) & Guarantee exercised & D & A & R & X & E & B & P & Pass\\\midrule
Initial completion and no-op repair & 01, 02 & Thm.~\ref{thm:safety}, Lem.~\ref{thm:local} & 8 & 8 & 0 & 0 & 0 & 0 & 1 & 2/2\\
Source edits, additions, deletions & 03, 04, 05, 06, 17 & Thm.~\ref{thm:kernel}, Lem.~\ref{thm:local} & 28 & 28 & 0 & 0 & 6 & 0 & 0 & 5/5\\
Shared-context changes and writers & 07, 08, 09, 10, 24 & Lem.~\ref{thm:local} & 18 & 18 & 0 & 0 & 6 & 1 & 0 & 5/5\\
Obsolete, late, duplicate replies & 11, 12, 13 & Thm.~\ref{thm:safety}(iv), \ref{thm:refinement} & 3 & 0 & 1 & 4 & 3 & 0 & 0 & 3/3\\
Budget exhaustion and blocking & 14, 15 & Thm.~\ref{thm:cost}, \ref{thm:termination} & 4 & 0 & 4 & 0 & 2 & 4 & 0 & 2/2\\
Contract change and early cutoff & 16, 23 & Lem.~\ref{thm:local}, Thm.~\ref{thm:reconfig} & 10 & 9 & 0 & 1 & 3 & 0 & 0 & 2/2\\
Joint conflict, rollback, stale check & 18, 19, 25 & Thm.~\ref{thm:safety}(i),(v) & 14 & 13 & 1 & 0 & 3 & 2 & 0 & 3/3\\
Reviewer and authority changes & 20, 21, 22 & Thm.~\ref{thm:reconfig} & 10 & 9 & 0 & 1 & 4 & 3 & 1 & 3/3\\
Untrusted ingress, observer loss & 26, 27 & Sec.~\ref{sec:extensions} & 4 & 4 & 0 & 0 & 0 & 2 & 1 & 2/2\\
Finite policy language (exhaustive) & 28 & Thm.~\ref{thm:table} & 0 & 0 & 0 & 0 & 0 & 0 & 0 & 1/1\\
\midrule\textbf{All scenarios} & 01--28 & & 99 & 89 & 6 & 6 & 27 & 12 & 3 & 28/28\\
\bottomrule\end{tabular*}
\end{table*}


\textbf{Answer to RQ3.} The reference implementation behaved as the semantics prescribes for all ten change classes, including the refusals, and an independent checker confirmed every recorded trace.

% ---------------------------------------------------------------------
\subsection{RQ4: Assurance and Remaining Obligations}\label{sec:rq4}
Table~\ref{tab:assurance} summarizes the mechanized and exhaustive evidence. The four TLAPS modules discharge \ProvedObligations{} obligations with the pinned toolchain (TLAPS 1.6.0-pre, build 5f5037f, Z3 4.8.9; the recorded logs come from this build, and the current upstream release archive no longer matches the pinned checksum, so we did not substitute it). The preserved first development run failed three obligations in two modules: proof steps that lacked explicit references to construction and type assumptions. The repaired modules split those steps and cite the assumptions explicitly, which changed the number of obligations that TLAPS generates (from 21 to 25). Counts are therefore comparable only within one version of a module, and the first-trial sources were not archived; every run now archives its sources. No assumption equivalent to the safety conclusion was added. For the finite policy language, exhaustive enumeration of all \PolicyTables{} typed tables confirmed that the validator accepts exactly one table per policy, and all \EventWords{} event words of length at most six obey the trace properties of Theorem~\ref{thm:table}.

% Generated by experiments/tse_paper.py; do not edit.
\begin{table}[!t]
\centering
\caption{Assurance evidence. For proofs, \emph{Final} counts proved/total obligations in the clean run and \emph{First trial} the preserved failed development run. Executable counts are exhaustive where stated; the other checks cover the recorded executions only.}\label{tab:assurance}
\scriptsize\setlength{\tabcolsep}{2.5pt}
\begin{tabular*}{\columnwidth}{@{\extracolsep{\fill}}>{\raggedright\arraybackslash}p{24mm}>{\raggedright\arraybackslash}p{28mm}rl@{}}\toprule
Artifact & Scope & Final & First trial\\\midrule
\multicolumn{4}{@{}l}{\textit{Mechanized TLAPS lemmas (strict mode, fingerprints erased)}}\\
{\ttfamily Budget} & Counter bound for one episode & 8/8 & 1 of 7 failed\\
{\ttfamily ControllerSafety} & Commit guard implies current, authorized, valid commit & 2/2 & 0 of 2 failed\\
{\ttfamily StructuralKernel} & Kernel coverage, unique producers, reference range & 9/9 & 2 of 6 failed\\
{\ttfamily DocumentPolicy} & Validator equivalence; archived-put conflict & 6/6 & 0 of 6 failed\\
\midrule\multicolumn{4}{@{}l}{\textit{Executable checks}}\\
Policy tables & All typed six-row tables, both policies & 8,192 & exhaustive\\
Event words & All words of length $\le 6$, strict table & 1,093 & exhaustive\\
Trace checker & Recorded scenario traces admitted & 28/28 & --\\
Trace checker & Deliberately corrupted traces rejected & 9/9 & --\\
Regression tests & Transaction-boundary unit tests & 14/14 & --\\
Reproduction & Clean re-runs of packaged archives matched (before the revision) & 2/2 & --\\
\bottomrule\end{tabular*}
\end{table}


Table~\ref{tab:tlc-ext} reports the extended model-checking campaign of this revision. It answers a limitation of the original submission, in which no multi-job instance with environmental changes had more than \TLCMaxStates{} states and all instances with three or more jobs had frozen inputs. The \TLCExtBaselines{} baselines include chains of up to \TLCExtMaxJobs{} jobs, a fan-out, and a diamond whose sink reads two prerequisites. Their inputs change up to twice, with budgets up to two, and they reach \TLCExtMaxStates{} distinct (\TLCExtMaxGenerated{} generated) states. Every baseline satisfies the full safety invariant and the refinement mapping of Theorem~\ref{thm:refinement}. With both obsolete-reply defenses removed, the refinement check fails, as it must. All \TLCExtRuns{} runs took \TLCExtSeconds\,s.

% Generated by experiments/revision_paper.py; do not edit.
\begin{table}[!t]
\centering
\caption{RQ4: extended bounded model checking (TLC 2026.11.19.180144 (rev: bc1b506), 303\,s for all 43 runs). $J$: jobs; $k$: attempt cap; $m$: environmental changes. Every baseline checks the full safety invariant \emph{and} the refinement mapping of Theorem~\ref{thm:refinement}. Mutant rows give distinct states and counterexample lengths for the three topologies.}\label{tab:tlc-ext}
\scriptsize\setlength{\tabcolsep}{3pt}\resizebox{\columnwidth}{!}{%
\begin{tabular}{@{}lrrrrrl@{}}\toprule
Configuration & $J$ & $k$ & $m$ & Distinct & Generated & Result\\\midrule
\multicolumn{7}{@{}l}{\textit{Baselines: \texttt{Safety} and refinement}}\\
One job & 1 & 2 & 2 & 4,107 & 24,095 & Pass\\
Chain & 2 & 2 & 1 & 1,563 & 10,073 & Pass\\
Chain & 2 & 1 & 2 & 7,740 & 42,520 & Pass\\
Chain & 3 & 1 & 1 & 634 & 3,343 & Pass\\
Chain & 3 & 2 & 1 & 6,735 & 54,814 & Pass\\
Fan-out & 3 & 1 & 1 & 1,327 & 7,856 & Pass\\
Fan-out & 3 & 1 & 2 & 64,932 & 465,720 & Pass\\
Diamond & 4 & 1 & 1 & 1,976 & 12,489 & Pass\\
Diamond & 4 & 2 & 1 & 42,623 & 440,879 & Pass\\
Diamond & 4 & 1 & 2 & 128,215 & 1,000,134 & Pass\\
Chain & 4 & 1 & 1 & 1,132 & 6,832 & Pass\\
Chain & 4 & 1 & 2 & 69,840 & 525,121 & Pass\\
Chain & 5 & 1 & 1 & 1,804 & 12,233 & Pass\\
Chain & 3 & 2 & 2 & 718,672 & 7,483,518 & Pass\\
\midrule\multicolumn{7}{@{}l}{\textit{Mutants on chains of 2 and 3 jobs and the 4-job diamond (shortest counterexample, states)}}\\
\multicolumn{4}{@{}l}{Cancellation removed} & \multicolumn{2}{r}{319/700/2,267} & Survives\textsuperscript{a}\\
\multicolumn{4}{@{}l}{Commit recheck removed} & \multicolumn{2}{r}{292/634/1,976} & Survives\textsuperscript{a}\\
\multicolumn{4}{@{}l}{Both obsolete defenses} & \multicolumn{2}{r}{321/535/513} & Viol.\ 4/4/4\\
\multicolumn{4}{@{}l}{Exhausted budget reset} & \multicolumn{2}{r}{102/114/122} & Viol.\ 5/5/5\\
\multicolumn{4}{@{}l}{Suspension ignored} & \multicolumn{2}{r}{173/206/322} & Viol.\ 3/3/3\\
\multicolumn{4}{@{}l}{Authority ignored} & \multicolumn{2}{r}{283/250/309} & Viol.\ 4/4/4\\
\multicolumn{4}{@{}l}{Revision check removed} & \multicolumn{2}{r}{161/268/857} & Viol.\ 7/9/11\\
\multicolumn{4}{@{}l}{Contract omitted} & \multicolumn{2}{r}{252/350/953} & Viol.\ 4/4/4\\
\multicolumn{4}{@{}l}{Selected context omitted} & \multicolumn{2}{r}{252/436/711} & Viol.\ 4/4/4\\
\midrule\multicolumn{7}{@{}l}{\textit{Refinement nonvacuity (no obsolete defenses)}}\\
One job & 1 & 1 & 1 & 96 & 367 & Violated\\
Chain & 2 & 1 & 1 & 120 & 424 & Violated\\
\bottomrule\end{tabular}}
\par\smallskip\raggedright\footnotesize\textsuperscript{a}The other obsolete-reply defense still blocks the stale commit.
\end{table}


The evidence remains deliberately uneven, and Table~\ref{tab:evidence} makes this explicit. The scoped lemmas used inside the safety and budget proofs are mechanized. The inductive invariant and the refinement mapping are model-checked on bounded instances but proved for all instances only on paper, and the history-counting argument across episodes is proved only on paper. The implementation-level ablation, the scenario campaign and the differential test cover larger and concurrent executions, but by testing rather than proof. All stages of the replication package, including the campaigns of this revision, were re-run on the revised code; the archived clean re-runs of the original submission predate the revision.

\textbf{Answer to RQ4.} The scoped lemmas, the extended bounded model checks, including a checked refinement mapping, and the exhaustive checks corroborate the paper proofs at every level we could mechanize. The parameterized invariant and refinement remain open mechanization targets, for the technical reasons given in Section~\ref{sec:artifact}.

% ---------------------------------------------------------------------
% ---------------------------------------------------------------------
\subsection{RQ5: Real LLM Proposers}\label{sec:rq5}
RQ1--RQ4 quantify over all proposals, so they say nothing about how often real models produce the situations the controller guards against. RQ5 runs the document-service workflow of Section~\ref{sec:llm-workflows} with \LlmModelsN{} open models (\LlmCallsTotal{} logged calls).

% Generated by experiments/revision_paper.py; do not edit.
\begin{table*}[!t]
\centering
\caption{RQ5: initial completion of four-service workspaces by real LLM proposers (budget $k=3$, validator feedback on retries, one feedback epoch after a barrier conflict). \emph{First-attempt acceptance}: share of first attempts committed per stage. \emph{Failed}: share of jobs that exhausted their budget (their dependents are blocked, an explicit verdict). \emph{Conflicts}: workspaces whose locally valid designs violated the global identifier convention at the barrier.}\label{tab:llm}
\footnotesize\setlength{\tabcolsep}{4pt}
\begin{tabular*}{\textwidth}{@{\extracolsep{\fill}}lrrrrrrrrr@{}}\toprule
& & \multicolumn{3}{c}{First-attempt acceptance (\%)} & & & & & \\\cmidrule(lr){3-5}
Model & Trials & Design & Code & Tests & Failed jobs (\%) & Conflicts & Published & Calls & kTokens/call\\\midrule
gemma2:9b & 6 & 0 & 0 & -- & 33 & 0 & 0 & 98 & 0.6\\
gpt-oss:20b & 6 & 12 & 100 & 67 & 3 & 0 & 4 & 104 & 1.7\\
llama3.1:8b & 6 & 0 & 4 & 0 & 29 & 0 & 0 & 138 & 0.7\\
mistral-small:24b & 6 & 0 & 58 & 83 & 1 & 0 & 5 & 113 & 0.8\\
phi4:14b & 6 & 0 & 100 & 38 & 6 & 0 & 3 & 117 & 0.8\\
qwen2.5-coder:7b & 6 & 0 & 12 & 0 & 32 & 0 & 0 & 123 & 0.7\\
qwen3-coder:30b & 6 & 0 & 92 & 25 & 4 & 0 & 4 & 122 & 0.7\\
\bottomrule\end{tabular*}
\end{table*}


\emph{Completion.} Table~\ref{tab:llm} reports \LlmTrials{} fresh four-service workspaces, six per model, each with twelve jobs. \LlmPublished{} workspaces (\LlmPublishedPct\%) reached a published snapshot. In the others, \LlmJobsFailedPct\% of all jobs exhausted their budget of $k=3$ attempts, and their dependents were blocked. Every unpublished workspace ended with an explicit verdict that named the failed jobs. No run looped, and no snapshot was published with a failed, stale or pending job. \LlmModelsNeverPublished{} of the \LlmModelsN{} models, the three with fewer than 10 billion parameters, published no workspace at all. Across all models, the most frequent implementation rejection was a violation of the stated interface (\texttt{read} must always return a pair), and the most frequent test rejection was a suite that never exercised the archived-document rule; the smaller models failed these checks again on their retries more often. The controller cannot make a weak model succeed. What it guarantees is that failure is reported, bounded and contained.

\emph{First attempts and feedback.} First-attempt acceptance of designs is close to zero, and the logs show why. \LlmDesignNameRejects{} of the \LlmDesignFirstRejects{} rejected first design attempts failed the check on the list of operation names. Of the \LlmDesignParsed{} first attempts we could parse, \LlmDesignUpdateNamed{} named the update operation \texttt{update}, whereas the validator requires \texttt{put}. That name appears in the clauses (``Updating (put) \dots'') but not in the design contract. The contract is therefore underspecified: this is a defect of our workflow, not of the models. The validator's diagnostic names the expected list, and the second attempt usually repaired it. We report the runs as executed rather than re-running with a corrected contract, because they illustrate two points. First, validator feedback recovered \LlmFeedbackRecoveredPct\% (95\% CI \LlmFeedbackRecoveredCI\%) of all rejected first attempts across stages. Second, first-attempt rates measure the clarity of a contract at least as much as the capability of a model.

\emph{Stale replies.} The stale trials are a controlled experiment, not a field observation. An implementation request is dispatched, clause R03 is flipped half a second later, and the reply is delivered after the flip. Of the 84 planned trials, 7 could not run because the design job had already failed. In the remaining \LlmStaleTrials{}, the commit guard discarded every late reply, as Theorem~\ref{thm:safety} requires. To see what accepting them would have meant, we evaluated each late reply against the validators of both questions. \LlmStaleOldValid{} replies passed the validator of the question they were asked, and \LlmStaleWrongPct\% of these (95\% CI \LlmStaleWrongCI\%) failed the validator of the current question. The other replies were invalid for both questions. Because the flip changes exactly the behavior the implementation must encode, this outcome is expected. It nonetheless quantifies a practical point: acceptance rules of the form ``the reply is valid for the question it was asked'', which is what a framework without request identities and current keys enforces, would have committed \LlmStaleOldValid{} wrong implementations.

% Generated by experiments/revision_paper.py; do not edit.
\begin{table}[!t]
\centering
\caption{RQ5: early cutoff under stochastic generation. After a reworded contract with the same meaning, how often was the regenerated value equal to the previous one? \emph{Raw}: identical response text; \emph{Value}: identical slot value (what early cutoff compares); \emph{Sem.}: same identifier type, operations and status codes; \emph{AST}: same structure-only normal form.}\label{tab:llm-cutoff}
\scriptsize\setlength{\tabcolsep}{1.6pt}
\begin{tabular*}{\columnwidth}{@{\extracolsep{\fill}}lrrrrrrrrr@{}}\toprule
& & \multicolumn{4}{c}{Design (JSON), \%} & \multicolumn{4}{c}{Implementation (code), \%}\\\cmidrule(lr){3-6}\cmidrule(lr){7-10}
Model & $T$ & $n$ & Raw & Value & Sem. & $n$ & Raw & Value & AST\\\midrule
gpt-oss:20b & 0.7 & 8 & 12 & 62 & 100 & 8 & 0 & 0 & 0\\
llama3.1:8b & 0.7 & 1 & 0 & 100 & 100 & 1 & 0 & 0 & 0\\
mistral-small:24b & 0.0 & 8 & 50 & 100 & 100 & 8 & 50 & 50 & 50\\
mistral-small:24b & 0.7 & 8 & 62 & 100 & 100 & 8 & 25 & 25 & 38\\
phi4:14b & 0.0 & 8 & 100 & 100 & 100 & 8 & 0 & 0 & 0\\
phi4:14b & 0.7 & 5 & 80 & 100 & 100 & 5 & 0 & 0 & 0\\
qwen3-coder:30b & 0.0 & 8 & 0 & 100 & 100 & 8 & 12 & 12 & 12\\
qwen3-coder:30b & 0.7 & 8 & 25 & 100 & 100 & 8 & 12 & 12 & 12\\
\bottomrule\end{tabular*}
\end{table}


\emph{Early cutoff.} Table~\ref{tab:llm-cutoff} replaces a contract by a reworded instruction with the same meaning, so a perfect proposer would regenerate an equal value. Cutoff trials need a successful initial repair, so only models that completed one appear. For the design slot, whose value is canonical JSON, the regenerated value equaled the previous one in \LlmCutoffDesignValuePct\% of trials, although the raw response text was identical in only \LlmCutoffDesignRawPct\%. Canonicalization absorbs differences in formatting and key order, which is the case for narrowly typed slots. For free-form implementations, value equality dropped to \LlmCutoffImplValuePct\% (AST equality \LlmCutoffImplAstPct\%), and most regenerations propagated to the verification job.

% Generated by experiments/revision_paper.py; do not edit.
\begin{table}[!t]
\centering
\caption{RQ5: cost of one stream of eight changes (two clause edits, three context changes of which one unselected and one equal-content, two reworded contracts) on a four-service workspace with real LLMs, after the initial repair. \emph{Publishable}: changes after which the repaired workspace passed the barrier.}\label{tab:llm-cost}
\scriptsize\setlength{\tabcolsep}{2pt}
\begin{tabular*}{\columnwidth}{@{\extracolsep{\fill}}llrrrr@{}}\toprule
Model & Strategy & Requests & kTokens & Seconds & Publishable\\\midrule
mistral-small:24b & Exact (ours) & 34 & 31.4 & 128 & 5/8\\
mistral-small:24b & No early cutoff & 44 & 41.3 & 136 & 7/8\\
mistral-small:24b & Whole workspace & 117 & 202.5 & 288 & 3/8\\
mistral-small:24b & Full regeneration & 159 & 136.5 & 419 & 7/8\\
qwen2.5-coder:7b & Exact (ours) & 50 & 40.1 & 116 & 0/8\\
qwen2.5-coder:7b & No early cutoff & 50 & 41.4 & 106 & 0/8\\
qwen2.5-coder:7b & Whole workspace & 94 & 134.4 & 116 & 0/8\\
qwen2.5-coder:7b & Full regeneration & 170 & 118.0 & 227 & 0/8\\
qwen3-coder:30b & Exact (ours) & 37 & 28.0 & 117 & 5/8\\
qwen3-coder:30b & No early cutoff & 43 & 33.3 & 143 & 5/8\\
qwen3-coder:30b & Whole workspace & 132 & 208.5 & 454 & 1/8\\
qwen3-coder:30b & Full regeneration & 154 & 107.1 & 416 & 7/8\\
\bottomrule\end{tabular*}
\end{table}


\emph{Cost and outcome.} Table~\ref{tab:llm-cost} applies one fixed stream of eight changes under four strategies with three models. Over the three models, exact keys needed \LlmCostExactReq{} requests against \LlmCostFullReq{} for full regeneration (\LlmCostRatio$\times$) and \LlmCostTokenRatio$\times$ fewer tokens. The whole-workspace strategy sent the most tokens, because every prompt carried the whole workspace, and it produced the fewest publishable snapshots. Exact keys did \emph{not}, however, dominate on outcome. For \texttt{mistral-small} they yielded a publishable snapshot after 5 of 8 changes, against 7 of 8 for full regeneration. The trace explains why. One verification job exhausted its budget at the first change, and its failed verdict persisted through the next two changes, which did not alter its observation, until a contract change opened a new episode. Full regeneration retries every job after every change and so re-samples a failed job by chance. Exact keys retain the verdict, as the semantics prescribes (Theorem~\ref{thm:termination}): Failed is terminal for an unchanged episode. This is a deliberate trade-off, not an accident. A deployment that wants more attempts can open a new episode explicitly, for example by changing the job's generation settings, which are part of its contract. It then pays for that job only, not for the whole workspace.

\textbf{Answer to RQ5.} With real models the guard was exercised as intended. Every late reply was discarded, and a large majority of those that looked valid were wrong for the current question. Validator feedback repaired most rejected attempts. Exact keys cut requests by \LlmCostRatio$\times$ against full regeneration. Early cutoff fired for canonical structured slots but rarely for free-form code. Two results qualify the picture. Smaller models often failed outright, and the controller can only make such failure explicit. And because failure verdicts are remembered, exact keys can publish less often than blind regeneration unless retries are requested explicitly. The barrier conflict that motivates the global predicate (F2) did not occur with these models: every model chose the same identifier type for all services (\LlmBarrierConflicts{} conflicts), so the evidence for the barrier remains RQ1 and the model checker.

% ---------------------------------------------------------------------
\subsection{RQ6: Validator Strength on Standard Benchmarks}\label{sec:rq6}
RQ6 asks what coordination guarantees mean for correctness when the validator is a typical test suite. It uses the \BenchTasks{} tasks of HumanEval+ and MBPP+ with \BenchModelsN{} models (\BenchCalls{} calls). The benchmark's base tests serve as the declared validator, and its plus tests serve only as hidden ground truth.

% Generated by experiments/revision_paper.py; do not edit.
\begin{table*}[!t]
\centering
\caption{RQ6: coordinated code generation on HumanEval+ (HE+) and MBPP+ with budget $k=3$ and the benchmark's base tests as the declared validator; plus tests (up to 200 inputs per task) are used only as ground truth. \emph{First correct}: first proposal passes base and plus tests (no validator, no retry). \emph{Accepted}: a value became usable. \emph{Accepted \& correct}: it also passes the plus tests. \emph{Wrong among accepted}: validator-relative correctness gap (Corollary~\ref{cor:domain}). Right block: after a reworded code contract with the same meaning, share of regenerated values that are identical as text (\emph{Raw}), identical as slot values (\emph{Value}; early cutoff fires), identical in structure (\emph{AST}), or behaviourally equal on all plus inputs (\emph{Beh.}).}\label{tab:bench}
\footnotesize\setlength{\tabcolsep}{3.5pt}
\begin{tabular*}{\textwidth}{@{\extracolsep{\fill}}llrrrrrrrrr@{}}\toprule
& & & \multicolumn{4}{c}{Correctness (\%)} & \multicolumn{4}{c}{Regeneration equal to previous value (\%)}\\\cmidrule(lr){4-7}\cmidrule(lr){8-11}
Model & Bench. & Tasks & First correct & Accepted & Accepted \& correct & Wrong among accepted & Raw & Value & AST & Beh.\\\midrule
gemma2:9b & HE+ & 164 & 61.6 & 73.2 & 65.9 & 10.0 & 61 & 63 & 75 & 97\\
gemma2:9b & MBPP+ & 378 & 61.1 & 76.2 & 64.6 & 15.3 & 62 & 71 & 75 & 96\\
llama3.1:8b & HE+ & 164 & 61.0 & 75.0 & 68.3 & 8.9 & 38 & 38 & 55 & 97\\
llama3.1:8b & MBPP+ & 378 & 64.3 & 78.3 & 68.5 & 12.5 & 11 & 13 & 59 & 91\\
mistral-small:24b & HE+ & 164 & 81.1 & 90.2 & 84.8 & 6.1 & 33 & 33 & 60 & 97\\
mistral-small:24b & MBPP+ & 378 & 67.7 & 86.2 & 71.2 & 17.5 & 6 & 6 & 11 & 95\\
phi4:14b & HE+ & 164 & 81.7 & 91.5 & 85.4 & 6.7 & 18 & 18 & 43 & 95\\
phi4:14b & MBPP+ & 378 & 71.7 & 88.6 & 73.8 & 16.7 & 6 & 6 & 20 & 91\\
qwen2.5-coder:7b & HE+ & 164 & 81.7 & 89.6 & 85.4 & 4.8 & 50 & 50 & 60 & 97\\
qwen2.5-coder:7b & MBPP+ & 378 & 71.2 & 87.3 & 74.1 & 15.2 & 21 & 21 & 42 & 93\\
\bottomrule\end{tabular*}
\end{table*}


\emph{Correctness.} Table~\ref{tab:bench} reports the results. Over all models and tasks, a value was accepted for \BenchAcceptedPct\% of tasks. \BenchAcceptedWrongPct\% of the accepted values (95\% CI \BenchAcceptedWrongCI\%) failed the plus tests: \BenchAcceptedWrongHE\% on HumanEval+ and \BenchAcceptedWrongMBPP\% on MBPP+, whose base tests are known to be weaker~\cite{evalplus}. The gap is the one Corollary~\ref{cor:domain} predicts. The controller guarantees that an accepted value passed the declared validator for the current question, and the plus tests show how much that validator leaves unchecked. Coordination neither creates this gap nor hides it. Validator-guided retry raised the share of tasks with a correct final value from \BenchFirstCorrectPct\% (first proposal, no retry) to \BenchFinalCorrectPct\%, a gain of \BenchRetryGainPts{} percentage points. Of the first proposals that passed the base tests, \BenchFirstBaseWrongPct\% still failed the plus tests.

\emph{Regeneration under a reworded contract.} After every code contract was replaced by a reworded instruction with the same meaning, \BenchRegenN{} tasks were regenerated with a usable value. The new value was identical as text in \BenchByteEqualPct\% of them, equal as a slot value in \BenchValueEqualPct\% (95\% CI \BenchValueEqualCI\%) and equal in AST structure in \BenchAstEqualPct\%. It was behaviorally equal on all plus inputs in \BenchBehEqualPct\% (95\% CI \BenchBehEqualCI\%). Early cutoff fires only on value equality, so the test job was regenerated for \BenchTestsRegenPct\% of tasks, although almost all regenerated functions behaved identically. The gap between value and behavioral equality is the saving that a coarser equality could recover, at the price of false hits (Section~\ref{sec:discussion}).

\textbf{Answer to RQ6.} On standard benchmarks, correctness of the coordinated output was bounded by the declared validator, exactly as Corollary~\ref{cor:domain} states. With the base tests as validator, about one accepted value in eight was wrong, and the rate depended on the benchmark's test strength, not on coordination. For free-form code, stochastic regeneration made exact early cutoff fire in a minority of cases, even though behavior was almost always preserved.

% ---------------------------------------------------------------------
\subsection{Discussion: What the Evidence Adds}\label{sec:discussion}
Each guarantee has a well-known ancestor (Section~\ref{sec:related}). For each one we ask what our evidence shows beyond that ancestor, and where it shows nothing.

\emph{Currentness of content, not only of reads (Theorem~\ref{thm:safety}; RQ1, RQ5).} OCC validates at commit that the items a transaction read are unchanged~\cite{occ}. Recent transactional proposals for LLM agents carry this check, or compensation, over to agent actions~\cite{chang2025sagallm,lyu2026coagent,yang2026concurrency}. Such a check corresponds to our request-identity defense over the data a job reads. RQ1 shows that it is not sufficient for generated content. When the contract or the selected context was left out of the key, while all source reads were still checked, stale values became usable in \ContractStalePct\% and \ContextStalePct\% of runs. A generated artifact becomes stale when anything that shaped the question changes: the instruction, the validator, the selected evidence or the owner. RQ5 shows that this matters with real models. Of \LlmStaleTrials{} replies that arrived after their requirement changed, \LlmStaleOldValid{} would have passed the old validator, and \LlmStaleWrongPct\% of those were wrong under the new one. A late reply is therefore not merely outdated; it is usually incorrect. A design that asks the producing model whether a conflicting change invalidates its output~\cite{lyu2026coagent} would delegate this judgment to the component whose replies were wrong.

\emph{Early cutoff under nondeterminism (Theorems~\ref{thm:local} and~\ref{thm:cost}; RQ2, RQ5, RQ6).} Early cutoff is standard in build systems~\cite{mitchell2012shake,mokhov,erdweg2015pluto,hammer2014adapton}, where deterministic tasks make it pay off reliably. With stochastic producers, equality-based cutoff stays sound, because an equal value is the same evidence, but its benefit becomes an empirical question. Neither the build-system literature nor studies of LLM nondeterminism~\cite{ouyang2025nondeterminism} answer it. RQ2 isolates the structural ceiling: with deterministic proposers, cutoff reduced a cosmetic edit of a shared model from \CoupledCosmeticNoCut{} to \CoupledCosmeticExact{} requests. RQ5 and RQ6 measure how much of that ceiling survives real generation. Narrow, canonical slots kept most of it: regenerated designs were value-equal in \LlmCutoffDesignValuePct\% of trials, but raw text was identical in only \LlmCutoffDesignRawPct\%. Free-form code lost most of it. Only \BenchValueEqualPct\% of regenerated benchmark functions were value-equal, although \BenchBehEqualPct\% were behaviorally equal on all plus inputs. The gap between these two numbers is the saving that a semantic cache~\cite{bang2023gptcache,schroeder2026vcache} or a contract-level normal form could recover, at the price of false hits, which exact keys exclude by construction. Slot typing is thus a design parameter, and our data quantify its effect.

\emph{Bounded repair, not bounded retries (Theorems~\ref{thm:single} and~\ref{thm:termination}, Proposition~\ref{prop:eager}).} Feedback-driven retry is a standard ingredient of LLM code generation~\cite{madaan2023selfrefine,shinn2023reflexion,chen2024selfdebug,huang2023agentcoder,schafer2024testpilot}, usually with a per-run iteration cap. Our budget belongs to a dependency episode, which is why the terminal outcome holds under arbitrary interleavings of edits. The upstream gate is why a change starts at most one episode per affected job. The scheduling experiment shows that the gate is not an implementation detail. Without it, the same lattices needed up to \SchedEagerReqRatioMax$\times$ the requests and gained no latency. RQ5 shows the budget at work with real models: retries with validator feedback recovered \LlmFeedbackRecoveredPct\% of failed first attempts. The remaining failures (\LlmJobsFailedPct\% of jobs) ended in explicit verdicts that blocked their dependents, never in silent acceptance or unbounded retry. Section~\ref{sec:rq5} discusses what this means for weaker models, for which most workspaces ended blocked.

\emph{Locally valid is not jointly valid (global predicate, P7; Proposition~\ref{prop:epochs}).} Benchmarks of parallel coding agents report patches that pass alone and fail together~\cite{xia2026stale,khatua2026cooperbench}. This is the LLM counterpart of the result from model consistency that pairwise consistency does not imply global consistency~\cite{stevens}. Our barrier checks the global predicate before publication, and RQ1 shows that removing it publishes inconsistent snapshots (\GlobalPublishPct\% of runs). With real models, however, the barrier never fired in our workflow (\LlmBarrierConflicts{} conflicts in \LlmTrials{} trials): every model chose the same identifier type for all services. Our evidence for the barrier and for the feedback-epoch bound is therefore the adversarial campaign and the model checker, not RQ5. How often teams of real agents produce jointly inconsistent artifacts depends on the workflow, and the studies cited above suggest that richer interfaces make it common.

\emph{Validator-relative correctness (Corollary~\ref{cor:domain}; RQ6).} It is known that weak tests accept wrong programs~\cite{evalplus}, and multi-agent code generators already pair coders with testers~\cite{dong2024selfcollab,huang2023agentcoder,chen2023codet}. RQ6 adds a clean attribution inside a coordinated workflow. Coordination neither created nor hid incorrectness: \BenchAcceptedWrongPct\% of accepted values failed the hidden plus tests, a figure set by the strength of the base tests that served as validator. The corollary states this dependence, and the experiment measures its size.

\emph{Verified once, not per task (Theorems~\ref{thm:refinement} and~\ref{thm:reconfig}; RQ4).} MAS model checkers verify agents~\cite{lomuscio2017mcmas,dennis2012ajpf}, and recent work verifies LLM-agent protocols or workflows per task~\cite{xia2026tracefix,wang2026lean4agent}. Our verification object is the task-independent kernel, and this answers the question of who writes the specifications. A user of the kernel writes jobs, read sets and validators, which are executable checks, as in RQ5 and RQ6, not TLA$^+$. The model and proofs are written once by the kernel's developers, as for verified infrastructure in industry~\cite{newcombe2015aws,howard2025smartcasual}. Trace validation~\cite{cirstea2024trace} is the established path for keeping the implementation aligned with the model. Studies of LLM-written TLA$^+$ report low semantic correctness~\cite{bisharat2026tlallm}, which supports this division of labor rather than asking users or models to specify each workflow.


% ---------------------------------------------------------------------
\subsection{Implications for Builders of Agent Frameworks}\label{sec:implications}
The results suggest concrete rules for orchestrators of LLM agents, including graph-based ones such as LangGraph~\cite{langgraph}. Each rule can be implemented as a commit policy on top of typed shared state.
\begin{enumerate}
\item \emph{Key every generated artifact by its exact observation, including the prompt, validator version and owner.} RQ1 shows that omitting the contract or the selected context exposes stale artifacts, and RQ2 that content-based keys avoid most regeneration.
\item \emph{Give every request an identity, and recheck currentness at commit time.} Either defense alone prevents stale commits, but they serve different purposes, so keep both: the recheck makes acceptance safe, and cancellation frees the job for re-dispatch. RQ5 shows that the late replies this rejects are frequently wrong, not merely outdated.
\item \emph{Separate ``accepted'' from ``usable'', and gate consumers on usability.} Removing the upstream gate was the most frequently detected ablation (RQ1), and the gate also makes repair cost independent of scheduling (RQ2).
\item \emph{Make completion a global predicate, recomputed after every change of team or rules}, and admit repair loops only as bounded feedback epochs.
\item \emph{Budget per dependency episode, not per scheduler pass}, and feed validator diagnostics back to the next attempt (RQ5).
\item \emph{Type slots narrowly.} Early cutoff fires only when regenerated values are equal; narrowly typed, canonicalized slots make that common, and free-form slots make it rare (RQ5, RQ6).
\item \emph{Treat validator strength as the bound on correctness.} Coordination does not create correctness that the validators do not check (RQ6).
\end{enumerate}
None of these rules depends on how good the underlying model is.
